% Options for packages loaded elsewhere
% Options for packages loaded elsewhere
\PassOptionsToPackage{unicode}{hyperref}
\PassOptionsToPackage{hyphens}{url}
\PassOptionsToPackage{dvipsnames,svgnames,x11names}{xcolor}
%
\documentclass[
  12pt,
  letterpaper,
]{article}
\usepackage{xcolor}
\usepackage[margin=1in]{geometry}
\usepackage{amsmath,amssymb}
\setcounter{secnumdepth}{-\maxdimen} % remove section numbering
\usepackage{iftex}
\ifPDFTeX
  \usepackage[T1]{fontenc}
  \usepackage[utf8]{inputenc}
  \usepackage{textcomp} % provide euro and other symbols
\else % if luatex or xetex
  \usepackage{unicode-math} % this also loads fontspec
  \defaultfontfeatures{Scale=MatchLowercase}
  \defaultfontfeatures[\rmfamily]{Ligatures=TeX,Scale=1}
\fi
\usepackage{lmodern}
\ifPDFTeX\else
  % xetex/luatex font selection
  \setmainfont[]{Times New Roman}
\fi
% Use upquote if available, for straight quotes in verbatim environments
\IfFileExists{upquote.sty}{\usepackage{upquote}}{}
\IfFileExists{microtype.sty}{% use microtype if available
  \usepackage[]{microtype}
  \UseMicrotypeSet[protrusion]{basicmath} % disable protrusion for tt fonts
}{}
\usepackage{setspace}
\makeatletter
\@ifundefined{KOMAClassName}{% if non-KOMA class
  \IfFileExists{parskip.sty}{%
    \usepackage{parskip}
  }{% else
    \setlength{\parindent}{0pt}
    \setlength{\parskip}{6pt plus 2pt minus 1pt}}
}{% if KOMA class
  \KOMAoptions{parskip=half}}
\makeatother
% Make \paragraph and \subparagraph free-standing
\makeatletter
\ifx\paragraph\undefined\else
  \let\oldparagraph\paragraph
  \renewcommand{\paragraph}{
    \@ifstar
      \xxxParagraphStar
      \xxxParagraphNoStar
  }
  \newcommand{\xxxParagraphStar}[1]{\oldparagraph*{#1}\mbox{}}
  \newcommand{\xxxParagraphNoStar}[1]{\oldparagraph{#1}\mbox{}}
\fi
\ifx\subparagraph\undefined\else
  \let\oldsubparagraph\subparagraph
  \renewcommand{\subparagraph}{
    \@ifstar
      \xxxSubParagraphStar
      \xxxSubParagraphNoStar
  }
  \newcommand{\xxxSubParagraphStar}[1]{\oldsubparagraph*{#1}\mbox{}}
  \newcommand{\xxxSubParagraphNoStar}[1]{\oldsubparagraph{#1}\mbox{}}
\fi
\makeatother


\usepackage{longtable,booktabs,array}
\usepackage{calc} % for calculating minipage widths
% Correct order of tables after \paragraph or \subparagraph
\usepackage{etoolbox}
\makeatletter
\patchcmd\longtable{\par}{\if@noskipsec\mbox{}\fi\par}{}{}
\makeatother
% Allow footnotes in longtable head/foot
\IfFileExists{footnotehyper.sty}{\usepackage{footnotehyper}}{\usepackage{footnote}}
\makesavenoteenv{longtable}
\usepackage{graphicx}
\makeatletter
\newsavebox\pandoc@box
\newcommand*\pandocbounded[1]{% scales image to fit in text height/width
  \sbox\pandoc@box{#1}%
  \Gscale@div\@tempa{\textheight}{\dimexpr\ht\pandoc@box+\dp\pandoc@box\relax}%
  \Gscale@div\@tempb{\linewidth}{\wd\pandoc@box}%
  \ifdim\@tempb\p@<\@tempa\p@\let\@tempa\@tempb\fi% select the smaller of both
  \ifdim\@tempa\p@<\p@\scalebox{\@tempa}{\usebox\pandoc@box}%
  \else\usebox{\pandoc@box}%
  \fi%
}
% Set default figure placement to htbp
\def\fps@figure{htbp}
\makeatother





\setlength{\emergencystretch}{3em} % prevent overfull lines

\providecommand{\tightlist}{%
  \setlength{\itemsep}{0pt}\setlength{\parskip}{0pt}}



 


\usepackage{booktabs}
\usepackage{longtable}
\usepackage{float}
\usepackage{setspace}
\doublespacing
% PAR heading styles
\usepackage{titlesec}
% Principal subheads: 14pt bold
\titleformat{\section}{\normalfont\fontsize{14}{16}\bfseries}{\thesection}{1em}{}
% Secondary subheads: 12pt bold
\titleformat{\subsection}{\normalfont\fontsize{12}{14}\bfseries}{\thesubsection}{1em}{}
% Sub-subheads: 12pt bold italic, run-in
\titleformat{\subsubsection}[runin]{\normalfont\fontsize{12}{14}\bfseries\itshape}{\thesubsubsection}{1em}{}[.]
% Abstract formatting
\renewenvironment{abstract}{\par\bigskip\noindent\textbf{Abstract}\par\smallskip\noindent}{\par\bigskip}
% Remove section numbers
\setcounter{secnumdepth}{0}
\makeatletter
\@ifpackageloaded{caption}{}{\usepackage{caption}}
\AtBeginDocument{%
\ifdefined\contentsname
  \renewcommand*\contentsname{Table of contents}
\else
  \newcommand\contentsname{Table of contents}
\fi
\ifdefined\listfigurename
  \renewcommand*\listfigurename{List of Figures}
\else
  \newcommand\listfigurename{List of Figures}
\fi
\ifdefined\listtablename
  \renewcommand*\listtablename{List of Tables}
\else
  \newcommand\listtablename{List of Tables}
\fi
\ifdefined\figurename
  \renewcommand*\figurename{Figure}
\else
  \newcommand\figurename{Figure}
\fi
\ifdefined\tablename
  \renewcommand*\tablename{Table}
\else
  \newcommand\tablename{Table}
\fi
}
\@ifpackageloaded{float}{}{\usepackage{float}}
\floatstyle{ruled}
\@ifundefined{c@chapter}{\newfloat{codelisting}{h}{lop}}{\newfloat{codelisting}{h}{lop}[chapter]}
\floatname{codelisting}{Listing}
\newcommand*\listoflistings{\listof{codelisting}{List of Listings}}
\makeatother
\makeatletter
\makeatother
\makeatletter
\@ifpackageloaded{caption}{}{\usepackage{caption}}
\@ifpackageloaded{subcaption}{}{\usepackage{subcaption}}
\makeatother
\usepackage{bookmark}
\IfFileExists{xurl.sty}{\usepackage{xurl}}{} % add URL line breaks if available
\urlstyle{same}
\hypersetup{
  pdftitle={Appendix A: Data Description},
  colorlinks=true,
  linkcolor={blue},
  filecolor={Maroon},
  citecolor={Blue},
  urlcolor={Blue},
  pdfcreator={LaTeX via pandoc}}


\title{Appendix A: Data Description}
\author{}
\date{}
\begin{document}
\maketitle


\setstretch{2}
\subsection{A.1 Data Sources}\label{a.1-data-sources}

\subsubsection{Quarterly Performance Reports
(QPR)}\label{quarterly-performance-reports-qpr}

The QPR data is extracted from HUD's DRGR (Disaster Recovery Grant
Reporting) system. This administrative data captures:

\begin{itemize}
\tightlist
\item
  Fund allocations by appropriation
\item
  Obligation, disbursement, and expenditure amounts
\item
  Activity types (housing, infrastructure, economic development,
  planning, administration)
\item
  Quarterly reporting periods (2003-2023)
\end{itemize}

The raw dataset contains 130,605 quarterly observations across 78 unique
grantees and 206 grant allocations.

\subsubsection{Population Data}\label{population-data}

Population figures are obtained from:

\begin{itemize}
\tightlist
\item
  U.S. Census Bureau decennial census (2000, 2010, 2020)
\item
  American Community Survey estimates for intercensal years
\end{itemize}

\subsubsection{Disaster Characteristics}\label{disaster-characteristics}

Disaster type classifications derive from:

\begin{itemize}
\tightlist
\item
  FEMA declared disaster data
\item
  Primary hazard type (hurricane, wildfire, flood, other)
\item
  Disaster declaration date for era classification
\end{itemize}

\subsubsection{Experience Indicators}\label{experience-indicators}

Grantee experience measures are computed from:

\begin{itemize}
\tightlist
\item
  Count of prior CDBG-DR grants managed
\item
  Binary classification: Novice (first CDBG-DR grant) vs.~Experienced
  (2+ prior grants)
\end{itemize}

\subsection{A.2 Spending Velocity
Construction}\label{a.2-spending-velocity-construction}

\subsubsection{The Dynamic Denominator
Problem}\label{the-dynamic-denominator-problem}

Initial velocity calculations using period-specific denominators
produced computational artifacts. When the denominator (obligated
amount) changes across periods, apparent velocity changes reflect
administrative adjustments rather than actual spending acceleration:

\[
\text{Velocity}_t^{\text{raw}} = \frac{\text{Expended}_t}{\text{Obligated}_t} - \frac{\text{Expended}_{t-1}}{\text{Obligated}_{t-1}}
\]

This formula produces spurious velocity when obligations change. For
example, if \(\text{Obligated}\) doubles from quarter \(t-1\) to \(t\)
while \(\text{Expended}\) increases moderately, the ratio drops despite
actual spending increases.

\subsubsection{Standardized Velocity with Fixed
Denominators}\label{standardized-velocity-with-fixed-denominators}

The standardized approach uses the final obligated amount as a stable
denominator:

\[
\text{Velocity}_t^{\text{std}} = \frac{\text{Expended}_t}{\text{Obligated}_{\text{final}}} - \frac{\text{Expended}_{t-1}}{\text{Obligated}_{\text{final}}}
\]

This ensures that velocity changes reflect only changes in the numerator
(actual expenditure), not administrative adjustments to obligations.

\subsubsection{Impact of
Standardization}\label{impact-of-standardization}

\begin{longtable}[]{@{}llll@{}}
\toprule\noalign{}
Metric & Raw Velocity & Standardized Velocity & Reduction \\
\midrule\noalign{}
\endhead
\bottomrule\noalign{}
\endlastfoot
Std Dev & 0.089 & 0.028 & 68.5\% \\
Extreme values (\textgreater{} & 0.10 & ) & 0.60\% \\
IQR & 0.031 & 0.012 & 61.3\% \\
\end{longtable}

The standardized approach dramatically reduces measurement noise while
preserving meaningful velocity variation.

\subsubsection{Winsorization}\label{winsorization}

Standardized velocity values are winsorized at the 1st and 99th
percentiles to limit outlier influence:

\[
\text{Velocity}^{\text{winsorized}} = \begin{cases}
P_1 & \text{if } \text{Velocity} < P_1 \\
P_{99} & \text{if } \text{Velocity} > P_{99} \\
\text{Velocity} & \text{otherwise}
\end{cases}
\]

This bounds extreme observations while preserving the full distribution
shape.

\subsection{A.3 Phase-Specific
Velocity}\label{a.3-phase-specific-velocity}

Programs are divided into temporal terciles based on elapsed quarters:

\begin{itemize}
\tightlist
\item
  \textbf{Early phase}: Quarters 1 through \(\lfloor T/3 \rfloor\)
\item
  \textbf{Middle phase}: Quarters \(\lfloor T/3 \rfloor + 1\) through
  \(\lfloor 2T/3 \rfloor\)
\item
  \textbf{Late phase}: Quarters \(\lfloor 2T/3 \rfloor + 1\) through
  \(T\)
\end{itemize}

where \(T\) is the total number of observed quarters.

Mean velocity within each phase creates phase-specific measures:

\[
\text{Velocity}_{\text{phase}} = \frac{1}{N_{\text{phase}}} \sum_{t \in \text{phase}} \text{Velocity}_t
\]

\subsection{A.4 Sample
Characteristics}\label{a.4-sample-characteristics}

\subsubsection{Overall Sample
Composition}\label{overall-sample-composition}

\begin{longtable}[]{@{}lrr@{}}
\toprule\noalign{}
\textbf{Characteristic} & \textbf{Count} & \textbf{Percent} \\
\midrule\noalign{}
\endhead
\bottomrule\noalign{}
\endlastfoot
\textbf{Total Observations} & 152 & 100.0\% \\
Unique grantees & 78 & -- \\
Unique disaster events & 18 & -- \\
Unique grant allocations & 206 & -- \\
\end{longtable}

\subsubsection{Government Type
Distribution}\label{government-type-distribution}

\begin{longtable}[]{@{}lrr@{}}
\toprule\noalign{}
\textbf{Government Type} & \textbf{N} & \textbf{Percent} \\
\midrule\noalign{}
\endhead
\bottomrule\noalign{}
\endlastfoot
State governments & 37 & 24.3\% \\
Local governments & 40 & 26.3\% \\
Multiple/Mixed & 75 & 49.3\% \\
\textbf{Total} & \textbf{152} & \textbf{100.0\%} \\
\end{longtable}

\subsubsection{Disaster Type
Distribution}\label{disaster-type-distribution}

\begin{longtable}[]{@{}lrr@{}}
\toprule\noalign{}
\textbf{Disaster Type} & \textbf{N} & \textbf{Percent} \\
\midrule\noalign{}
\endhead
\bottomrule\noalign{}
\endlastfoot
Hurricane & 89 & 58.6\% \\
Flood & 27 & 17.8\% \\
Wildfire & 24 & 15.8\% \\
Other & 12 & 7.9\% \\
\textbf{Total} & \textbf{152} & \textbf{100.0\%} \\
\end{longtable}

\subsubsection{Experience Distribution}\label{experience-distribution}

\begin{longtable}[]{@{}lrr@{}}
\toprule\noalign{}
\textbf{Experience Level} & \textbf{N} & \textbf{Percent} \\
\midrule\noalign{}
\endhead
\bottomrule\noalign{}
\endlastfoot
Novice (first CDBG-DR) & 74 & 48.7\% \\
Experienced (2+ prior) & 78 & 51.3\% \\
\textbf{Total} & \textbf{152} & \textbf{100.0\%} \\
\end{longtable}

\subsubsection{Program Type
Distribution}\label{program-type-distribution}

\begin{longtable}[]{@{}lrr@{}}
\toprule\noalign{}
\textbf{Primary Program Type} & \textbf{N} & \textbf{Percent} \\
\midrule\noalign{}
\endhead
\bottomrule\noalign{}
\endlastfoot
Housing & 67 & 44.1\% \\
Infrastructure & 41 & 27.0\% \\
Administration & 34 & 22.4\% \\
Other & 10 & 6.6\% \\
\textbf{Total} & \textbf{152} & \textbf{100.0\%} \\
\end{longtable}

\subsubsection{Censoring Status}\label{censoring-status}

\begin{longtable}[]{@{}lrr@{}}
\toprule\noalign{}
\textbf{Status} & \textbf{N} & \textbf{Percent} \\
\midrule\noalign{}
\endhead
\bottomrule\noalign{}
\endlastfoot
Completed (reached 95\% expenditure) & 40 & 26.3\% \\
Right-censored (not yet complete) & 112 & 73.7\% \\
\textbf{Total} & \textbf{152} & \textbf{100.0\%} \\
\end{longtable}

The high censoring rate (73.7\%) motivates survival analysis methods.
Standard regression approaches that drop censored observations would
discard 112 of 152 observations, severely limiting statistical power and
introducing selection bias.

\subsection{A.5 Velocity Summary
Statistics}\label{a.5-velocity-summary-statistics}

\begin{longtable}[]{@{}
  >{\raggedright\arraybackslash}p{(\linewidth - 12\tabcolsep) * \real{0.2029}}
  >{\raggedleft\arraybackslash}p{(\linewidth - 12\tabcolsep) * \real{0.1014}}
  >{\raggedleft\arraybackslash}p{(\linewidth - 12\tabcolsep) * \real{0.1449}}
  >{\raggedleft\arraybackslash}p{(\linewidth - 12\tabcolsep) * \real{0.1159}}
  >{\raggedleft\arraybackslash}p{(\linewidth - 12\tabcolsep) * \real{0.1304}}
  >{\raggedleft\arraybackslash}p{(\linewidth - 12\tabcolsep) * \real{0.1739}}
  >{\raggedleft\arraybackslash}p{(\linewidth - 12\tabcolsep) * \real{0.1304}}@{}}
\toprule\noalign{}
\begin{minipage}[b]{\linewidth}\raggedright
\textbf{Variable}
\end{minipage} & \begin{minipage}[b]{\linewidth}\raggedleft
\textbf{N}
\end{minipage} & \begin{minipage}[b]{\linewidth}\raggedleft
\textbf{Mean}
\end{minipage} & \begin{minipage}[b]{\linewidth}\raggedleft
\textbf{SD}
\end{minipage} & \begin{minipage}[b]{\linewidth}\raggedleft
\textbf{Min}
\end{minipage} & \begin{minipage}[b]{\linewidth}\raggedleft
\textbf{Median}
\end{minipage} & \begin{minipage}[b]{\linewidth}\raggedleft
\textbf{Max}
\end{minipage} \\
\midrule\noalign{}
\endhead
\bottomrule\noalign{}
\endlastfoot
Overall Velocity (pp/quarter) & 152 & 2.14 & 1.83 & 0.12 & 1.67 &
9.84 \\
Early Phase Velocity & 152 & 2.89 & 2.41 & 0.08 & 2.21 & 12.67 \\
Middle Phase Velocity & 152 & 1.92 & 1.76 & 0.04 & 1.48 & 8.93 \\
Late Phase Velocity & 152 & 1.61 & 1.89 & 0.02 & 1.12 & 11.24 \\
Duration (quarters, completed only) & 40 & 68.2 & 28.4 & 18 & 63.5 &
142 \\
\end{longtable}

\textbf{Key observations:}

\begin{enumerate}
\def\labelenumi{\arabic{enumi}.}
\item
  \textbf{Overall velocity} averages 2.14 percentage points per quarter,
  indicating that the typical grantee expends approximately 2\% of their
  total allocation each quarter.
\item
  \textbf{Phase pattern}: Velocity is highest in early phases (2.89
  pp/quarter) and declines through middle (1.92) and late (1.61) phases,
  consistent with front-loaded activity followed by closeout delays.
\item
  \textbf{Substantial variation} exists across grantees (SD = 1.83 for
  overall velocity), providing analytic leverage for detecting
  velocity-completion associations.
\item
  \textbf{Completion timing}: Among the 40 completed programs, median
  time to 95\% completion is 63.5 quarters (approximately 16 years).
\end{enumerate}

\subsection{A.6 Data Quality}\label{a.6-data-quality}

Quality assurance flags track potential measurement issues:

\begin{longtable}[]{@{}lrr@{}}
\toprule\noalign{}
\textbf{QA Flag} & \textbf{N} & \textbf{Percent} \\
\midrule\noalign{}
\endhead
\bottomrule\noalign{}
\endlastfoot
Extreme velocity ( & v & \textgreater{} 10 pp/quarter) \\
Large obligation adjustment (\textgreater25\%) & 12 & 7.9\% \\
Negative expenditure adjustment & 8 & 5.3\% \\
\end{longtable}

Sensitivity analyses excluding flagged observations produce
substantively similar results (see Appendix C).

\subsection{A.7 Disaster Coverage}\label{a.7-disaster-coverage}

The sample covers 18 major disaster events from 2003 to 2023:

\textbf{Hurricane Events:} - Hurricane Katrina (2005): Gulf Coast states
- Superstorm Sandy (2012): Mid-Atlantic and Northeast - Hurricane Harvey
(2017): Texas coast - Hurricane Maria (2017): Puerto Rico and U.S.
Virgin Islands - Hurricane Irma (2017): Florida and Caribbean

\textbf{Wildfire Events:} - California Wildfires (2017-2020): Multiple
events - Colorado Fires (2012-2013) - Arizona Wildfires (2011)

\textbf{Flood Events:} - Louisiana Floods (2016) - Colorado Floods
(2013) - Multiple smaller flood events

\textbf{Other Events:} - Joplin Tornado (2011) - Oklahoma Tornadoes
(2013) - Various smaller disasters

The diversity of disaster types, geographic locations, and governmental
structures provides variation needed to test context-specific velocity
effects.




\end{document}
